\chapter{Những câu hỏi về thất bại}
\markboth{Những câu hỏi về thất bại}{Những câu hỏi về thất bại}

\section{Em sợ sai}
\begin{truyen}{Em sợ sai}{Classroom Talk}
\chuhoa{K}{hi được gọi lên bảng, một bạn nói: ``Thầy ơi, em chưa chắc. Em sợ sai.''}

Thầy đáp: ``Ở lớp học mà còn sợ sai, sau này em sẽ sợ cả cuộc đời.''

Bạn ấy cười gượng: ``Sai trước lớp quê lắm thầy.''

Thầy nói: ``Quê một lần để biết sửa còn hơn im luôn để sai mãi.''

Bạn hít sâu rồi bước lên bảng.

\textit{(A student fears making mistakes in public. The teacher reminds the student that school is the safest place to be wrong and learn.)}
\end{truyen}

\begin{baihoc}
Lớp học sinh ra không phải để giấu cái sai, mà để tập sửa cái sai.
\textit{(A classroom exists not to hide mistakes, but to practise correcting them.)}
\end{baihoc}

\begin{ghinhoanh}
\textbf{Short English to remember:}

School is the safe place to be wrong.

Correction is part of learning.
\end{ghinhoanh}

\begin{tuhocnhanh}
1. Em đang sợ sai ở môn nào nhất? \textit{(In which subject are you most afraid of being wrong?)}

2. Em có thể thử sai an toàn một lần ở đâu? \textit{(Where can you try being wrong safely once?)}
\end{tuhocnhanh}
\ngancach

\section{Trượt một bài kiểm tra coi như xong?}
\begin{truyen}{Trượt một bài kiểm tra coi như xong?}{Classroom Talk}
\chuhoa{S}{au một bài điểm thấp, bạn nói: ``Chắc em xong rồi thầy.''}

Thầy hỏi: ``Em mới trượt một bài hay trượt luôn cuộc đời?''

Bạn cười buồn: ``Dạ một bài.''

Thầy nói: ``Vậy thì gọi đúng tên nó thôi. Một bài hỏng không có quyền định nghĩa cả em.''

Bạn ấy ngồi xuống, mắt bớt đỏ hơn.

\textit{(After one poor test, a student feels finished. The teacher narrows the failure back to its true size: one event, not a whole life.)}
\end{truyen}

\begin{baihoc}
Gọi đúng kích thước của thất bại sẽ giúp mình không tuyệt vọng quá mức.
\textit{(Naming failure at its true size keeps us from exaggerated despair.)}
\end{baihoc}

\begin{ghinhoanh}
\textbf{Short English to remember:}

One bad test is one event.

It is not your whole story.
\end{ghinhoanh}

\begin{tuhocnhanh}
1. Em có đang phóng to một thất bại gần đây không? \textit{(Are you enlarging a recent failure?)}

2. Tên thật của vấn đề đó là gì? \textit{(What is the real name of that problem?)}
\end{tuhocnhanh}
\ngancach

\section{Em thất bại vì xui}
\begin{truyen}{Em thất bại vì xui}{Classroom Talk}
\chuhoa{C}{ó bạn than: ``Em xui lắm thầy, trúng ngay phần em không học.''}

Thầy hỏi: ``Vậy phần em học có đủ rộng chưa?''

Bạn im.

Thầy nói tiếp: ``May rủi có thật. Nhưng mình không thể giao hết tương lai cho may rủi bằng cách học thiếu.''

Bạn ấy thở dài, nhưng lần này là kiểu thở dài đã hiểu.

\textit{(A student blames bad luck for failure. The teacher accepts luck exists, but refuses to let poor preparation hide behind it.)}
\end{truyen}

\begin{baihoc}
Người chuẩn bị kỹ có thể vẫn gặp xui, nhưng ít bị xui đánh ngã hơn.
\textit{(Prepared people may still meet bad luck, but they are less easily knocked down by it.)}
\end{baihoc}

\begin{ghinhoanh}
\textbf{Short English to remember:}

Luck is real.

Preparation still matters more.
\end{ghinhoanh}

\begin{tuhocnhanh}
1. Em có đang dùng chữ ``xui'' để che phần chưa chuẩn bị không? \textit{(Are you using luck to hide weak preparation?)}

2. Em cần chuẩn bị rộng hơn ở đâu? \textit{(Where do you need broader preparation?)}
\end{tuhocnhanh}
\ngancach

\section{Bị cười một lần nên không dám phát biểu nữa}
\begin{truyen}{Bị cười một lần nên không dám phát biểu nữa}{Classroom Talk}
\chuhoa{M}{ột bạn kể: ``Lần trước em nói sai, cả lớp cười nên giờ em không dám nói nữa.''}

Thầy hỏi: ``Em sợ mình sai nữa hay sợ người ta cười nữa?''

Bạn đáp: ``Cả hai.''

Thầy nói: ``Nếu vì một lần bị cười mà im luôn, em đang đưa quyền sống của mình cho người khác.''

Bạn ấy yên lặng rất lâu.

\textit{(A student stops speaking after being laughed at once. The teacher warns that giving up because of ridicule hands too much power to others.)}
\end{truyen}

\begin{baihoc}
Người khác có thể cười mình một lúc. Đừng để họ quyết định mình im lặng cả năm.
\textit{(Others may laugh for a moment. Do not let them decide your silence for a whole year.)}
\end{baihoc}

\begin{ghinhoanh}
\textbf{Short English to remember:}

Do not hand your voice away.

One laugh should not silence a year.
\end{ghinhoanh}

\begin{tuhocnhanh}
1. Một lần xấu hổ nào đang giữ em lại? \textit{(What embarrassing moment is still holding you back?)}

2. Em có thể thử lên tiếng lại ở tình huống nhỏ nào? \textit{(In what small situation can you speak again?)}
\end{tuhocnhanh}
\ngancach

\section{Nếu em từng lười thì còn cứu kịp không?}
\begin{truyen}{Nếu em từng lười thì còn cứu kịp không?}{Classroom Talk}
\chuhoa{M}{ột bạn hỏi rất thật: ``Thầy ơi, trước giờ em lười quá. Giờ cứu còn kịp không?''}

Thầy đáp: ``Kịp, nếu em chịu bắt đầu ngay thay vì tiếc tiếp.''

Bạn hỏi: ``Bắt đầu từ đâu?''

Thầy nói: ``Từ bài hôm nay. Người đổi được cuộc đời không phải người tiếc quá khứ giỏi, mà là người chịu sửa hiện tại.''

Bạn ấy ghi câu đó vào mép vở.

\textit{(A student asks whether it is too late after being lazy for a long time. The teacher shifts attention from regret to immediate action.)}
\end{truyen}

\begin{baihoc}
Quá khứ chỉ nguy hiểm khi mình cứ ngồi đó nhìn nó.
\textit{(The past becomes dangerous only when we keep sitting there staring at it.)}
\end{baihoc}

\begin{ghinhoanh}
\textbf{Short English to remember:}

It is not too late if you begin now.

Repair starts in the present.
\end{ghinhoanh}

\begin{tuhocnhanh}
1. Em tiếc điều gì trong việc học của mình? \textit{(What do you regret in your learning?)}

2. Việc nhỏ nào em có thể sửa ngay hôm nay? \textit{(What small thing can you fix today?)}
\end{tuhocnhanh}
\ngancach

\section{Thất bại có đáng xấu hổ không?}
\begin{truyen}{Thất bại có đáng xấu hổ không?}{Classroom Talk}
\chuhoa{M}{ột bạn hỏi: ``Thầy ơi, thất bại có đáng xấu hổ không?''}

Thầy đáp: ``Tùy.''

Bạn ngạc nhiên: ``Tùy là sao ạ?''

Thầy nói: ``Nếu em làm hết sức mà vẫn ngã, đó là chuyện phải học. Còn nếu em biết mình cần làm mà cứ trốn mãi, nỗi xấu hổ nằm ở chỗ bỏ mặc bản thân chứ không nằm ở kết quả.''

Bạn ấy gật đầu, vẻ mặt bớt rối hơn.

\textit{(A student asks whether failure is shameful. The teacher locates shame not in losing itself, but in refusing responsibility.)}
\end{truyen}

\begin{baihoc}
Không phải thất bại nào cũng đáng xấu hổ. Nhưng bỏ mặc khả năng của mình thì rất đáng tiếc.
\textit{(Not every failure is shameful. But neglecting your own potential is deeply regrettable.)}
\end{baihoc}

\begin{ghinhoanh}
\textbf{Short English to remember:}

Failure is not always shame.

Abandoning yourself is worse.
\end{ghinhoanh}

\begin{tuhocnhanh}
1. Em đang buồn vì kết quả hay vì biết mình đã buông? \textit{(Are you upset by the result or by knowing you gave up?)}

2. Em muốn sửa phần nào? \textit{(What part do you want to repair?)}
\end{tuhocnhanh}
\ngancach

\section{Em làm lại rồi mà vẫn sai}
\begin{truyen}{Em làm lại rồi mà vẫn sai}{Classroom Talk}
\chuhoa{C}{ó bạn nản lòng: ``Em làm lại mấy lần rồi mà vẫn sai thầy.''}

Thầy hỏi: ``Em đang làm lại cách cũ, hay đang sửa cách làm?''

Bạn im.

Thầy nói tiếp: ``Lặp lại chưa chắc là tiến bộ. Làm lại chỉ có giá trị khi mình nhìn ra chỗ phải đổi.''

Bạn ấy gật đầu như hiểu vì sao mình mệt mà vẫn đứng yên.

\textit{(A student repeats the work but still fails. The teacher distinguishes repetition from correction.)}
\end{truyen}

\begin{baihoc}
Kiên trì rất quý, nhưng kiên trì sai cách có thể làm mình mệt mà không đi tới đâu.
\textit{(Persistence is valuable, but persistence in the wrong way can exhaust us without progress.)}
\end{baihoc}

\begin{ghinhoanh}
\textbf{Short English to remember:}

Repeating is not always improving.

Change the method, not just the effort.
\end{ghinhoanh}

\begin{tuhocnhanh}
1. Em đang lặp lại hay đang sửa thật? \textit{(Are you repeating or truly correcting?)}

2. Cách làm nào em cần đổi? \textit{(What method do you need to change?)}
\end{tuhocnhanh}
\ngancach

\section{Ngã nhiều quá có lì không nổi}
\begin{truyen}{Ngã nhiều quá có lì không nổi}{Classroom Talk}
\chuhoa{M}{ột bạn hỏi rất nhỏ: ``Ngã nhiều quá thì có lì không nổi nữa không thầy?''}

Thầy nhìn bạn ấy rồi đáp: ``Có lúc có.''

Bạn cúi đầu.

Thầy nói chậm: ``Người ta không phải máy. Có lúc mệt đến mức không muốn đứng lên nữa. Nên ngoài chuyện cố gắng, em còn phải học cách nghỉ đúng, nhờ đúng người, và đừng tự chịu một mình quá lâu.''

Câu trả lời làm lớp học im hoàn toàn.

\textit{(A student asks whether too many falls can break resilience. The teacher answers with compassion and introduces rest and support as part of recovery.)}
\end{truyen}

\begin{baihoc}
Mạnh mẽ không có nghĩa là gồng mãi một mình.
\textit{(Strength does not mean carrying everything alone forever.)}
\end{baihoc}

\begin{ghinhoanh}
\textbf{Short English to remember:}

You are not a machine.

Rest and support matter too.
\end{ghinhoanh}

\begin{tuhocnhanh}
1. Em đang mệt vì chuyện gì lặp lại quá nhiều? \textit{(What repeated struggle is exhausting you?)}

2. Em có thể nhờ ai đúng lúc? \textit{(Whom can you ask for help at the right time?)}
\end{tuhocnhanh}
\ngancach

\section{Thua bạn thân có đau hơn không?}
\begin{truyen}{Thua bạn thân có đau hơn không?}{Classroom Talk}
\chuhoa{C}{ó bạn hỏi nửa cười nửa buồn: ``Sao thua bạn thân thấy đau hơn thua người lạ vậy thầy?''}

Thầy đáp: ``Vì người gần mình chạm vào lòng tự trọng của mình rõ hơn.''

Bạn hỏi tiếp: ``Vậy phải làm sao để đỡ cay?''

Thầy nói: ``Học cách vui cho người ta mà không tự bỏ cuộc với mình. Người khác đi nhanh không bắt em phải ghét họ, chỉ nhắc em nhìn lại đường mình.''

Bạn ấy gật đầu, nhưng kiểu gật đầu còn đang tập.

\textit{(A student feels extra pain when losing to a close friend. The teacher explains how proximity intensifies ego and comparison.)}
\end{truyen}

\begin{baihoc}
Người gần mình thành công có thể làm mình đau, nhưng cũng có thể giúp mình lớn nếu biết nhìn cho đúng.
\textit{(The success of those close to us can hurt, but it can also help us grow if we see it properly.)}
\end{baihoc}

\begin{ghinhoanh}
\textbf{Short English to remember:}

Close comparison hurts more.

Use it to grow, not to hate.
\end{ghinhoanh}

\begin{tuhocnhanh}
1. Em có đang ghen với ai gần mình không? \textit{(Are you jealous of someone close to you?)}

2. Em muốn chuyển cảm giác đó thành điều gì? \textit{(What do you want to turn that feeling into?)}
\end{tuhocnhanh}
\ngancach

\section{Làm sao biết nên cố tiếp hay dừng lại?}
\begin{truyen}{Làm sao biết nên cố tiếp hay dừng lại?}{Classroom Talk}
\chuhoa{M}{ột bạn hỏi: ``Khi thất bại rồi, làm sao biết nên cố tiếp hay dừng lại để đổi hướng hả thầy?''}

Thầy đáp: ``Phải hỏi hai điều. Một là mình còn tin giá trị của hướng đó không. Hai là mình đang thiếu thời gian hay thiếu sự phù hợp thật.''

Bạn hỏi: ``Nếu chưa rõ thì sao?''

Thầy nói: ``Thì đừng quyết trong lúc mệt nhất. Nhiều quyết định tệ được đưa ra đúng lúc lòng mình đau nhất.''

Bạn ấy gật đầu rất mạnh ở câu cuối.

\textit{(A student asks when to keep going and when to change direction. The teacher gives a framework and warns against deciding from exhaustion.)}
\end{truyen}

\begin{baihoc}
Không phải lúc nào dừng lại cũng là yếu. Nhưng quyết định dừng nên đến từ tỉnh táo, không phải từ kiệt sức.
\textit{(Stopping is not always weakness. But the decision should come from clarity, not exhaustion.)}
\end{baihoc}

\begin{ghinhoanh}
\textbf{Short English to remember:}

Do not choose a path while exhausted.

Clarity matters before quitting.
\end{ghinhoanh}

\begin{tuhocnhanh}
1. Em đang muốn bỏ vì không hợp hay vì quá mệt? \textit{(Do you want to quit because it does not fit, or because you are too tired?)}

2. Em cần làm rõ điều gì trước khi quyết? \textit{(What do you need to clarify before deciding?)}
\end{tuhocnhanh}
\ngancach